\section{Proofs}
\label{app:proofs}

\begin{proof}[Proof of \cref{lem:cone}]
By induction on $t$. For $t=0$, $h_i^{(0)} = H(E(x_i))$ depends on $W_i^0(x) = (x_i)$ through a
function that does not depend on $i$. Assume the claim for $t$. The update for $h_i^{(t+1)}$ reads
$h_j^{(t)}$ and $e_j$ for $|j-i|\le r$. Each $h_j^{(t)}$ is a fixed function of
$W_j^{rt}(x) \subseteq W_i^{r(t+1)}(x)$ (padding positions contribute fixed values, which are also
determined by the window because the window records $\bot$), and each $e_j$ is a function of
$\tilde x_j\in W_i^{r}(x)$. Since $U$ is shared, $h_i^{(t+1)}$ is a fixed function of
$W_i^{r(t+1)}(x)$. The decoder is pointwise, which gives the statement about $\hat y$.
\end{proof}

\begin{proof}[Proof of \cref{thm:lb}]
Let $\rho = rt < \rho_i(x)$ (because $t < \lceil\rho_i/r\rceil$ implies $rt < \rho_i$). By minimality
of $\rho_i(x)$, position $i$ is not $\rho$-determined, so there is $(x', i')$ with
$W_{i'}^{\rho}(x') = W_i^{\rho}(x)$ and $y_{i'}(x')\ne y_i(x)$. By \cref{lem:cone},
$\hat y_{i'}^{(t)}(x') = \hat y_i^{(t)}(x)$, so this common prediction differs from at least one of the
two targets. Taking $i$ with $\rho_i(x) = d^*(x)$ gives the statement about $\mathrm{LB}(x)$.
\end{proof}

\begin{proof}[Proof of \cref{cor:local}]
Suppose position $i$ of $x$ commits at $t < \lceil\rho_i(x)/r\rceil$. Take the pair $(x', i')$ from
\cref{thm:lb}. The histories $h_i^{(0..t)}(x)$ and $h_{i'}^{(0..t)}(x')$ coincide by
\cref{lem:cone}, so position $i'$ of $x'$ raises its flag at the same step $t$ (it cannot have
raised it earlier, since $x$ would then have raised it earlier too) and commits the same prediction,
which is wrong for one of them. The input halts when all positions have committed, so it halts no
earlier than $\max_i\lceil\rho_i(x)/r\rceil = \mathrm{LB}(x)$.
\end{proof}

\begin{proof}[Proof of \cref{cor:length}]
At length $L$ the input with $a_j+b_j = 9$ for all $j<L$ has no stops, so $d^* = L+1$ by
\cref{prop:oracles}, and \cref{thm:lb} applies to the final carry position. Any $\tau(L)$ below
$\lceil (L+1)/r\rceil$ is therefore wrong on one member of the corresponding pair, and the same
argument applies to every input with $d^*(x) > r\tau(L)$. For uniformly random digits each position
is a non-stop with probability $1/10$, independently, so $d^*-1$ is the longest run of successes in
$L$ Bernoulli$(1/10)$ trials, which is $\log_{10} L + O_p(1)$ (the Erd\H{o}s--R\'enyi longest-run law).
\end{proof}

\begin{proof}[Proof of \cref{prop:oracles}]
\emph{Addition.} Output $i$ is $(a_i + b_i + c_i)\bmod 10$ for $i<L$ and $c_L$ for $i=L$, where $c_i$ is
the carry into $i$. If $s_i\ge 0$ is the nearest stop before $i$, then positions $s_i+1,\dots,i-1$
all have $a_j+b_j = 9$, so $c_i = c_{s_i+1} = \mathbb 1[a_{s_i}+b_{s_i}\ge 10]$. If there is no stop,
all positions $0,\dots,i-1$ sum to 9 and $c_i = c_0 = 0$.
\emph{Upper bound} ($\rho_i \le i - s_i$): the window of radius $i - s_i$ contains position $s_i$ (or the
boundary symbol $\bot$ at $-1$) and every position between, so it fixes $c_i$ and $(a_i,b_i)$, hence
$y_i$, for every $(x',i')$ sharing the window. A boundary symbol at relative offset $-(i-s_i)$ in
$x'$ forces $i' = i$, so the carry is 0 there too.
\emph{Lower bound} ($\rho_i \ge i - s_i$): take $\rho = i - s_i - 1$. The window contains positions
$i-\rho,\dots,i-1$, all summing to 9, and no boundary symbol on the left. Build $x'$ by
taking the window content, placing it at $i' = \rho+1$, and putting one digit pair at
position $0$ of $x'$ whose sum is $\ge 10$ if $c_i(x) = 0$ and $<9$ if $c_i(x)=1$ (the right-hand
part of the window is copied unchanged). Then $c_{i'}(x') = 1 - c_i(x)$, so $y_{i'}(x')\ne y_i(x)$.
\emph{Parity} is identical with ``delimiter'' for ``stop'' and the segment prefix parity for the carry:
flipping the bit just outside the window changes the parity.
\end{proof}

\begin{proof}[Proof of \cref{prop:automaton}]
\emph{Addition.} Each position stores its pair sum class and a carry-in $c_i\in\{0,1,?\}$, initially
$?$. Per step, position $i$ reads its left neighbour's carry-out, which is
$\mathbb 1[a_{i-1}+b_{i-1}\ge 10]$ if $i-1$ is a stop, $c_{i-1}$ if $a_{i-1}+b_{i-1}=9$, and $0$ if
the neighbour is the boundary $\bot$; it sets $c_i$ to that value if $c_i = ?$ and the value is
known. This update reads only radius-1 neighbours and is the same at every position and length.
By induction, $c_i$ becomes known exactly at step $i - s_i = \rho_i$ (step 1 at a stop's right
neighbour or at position 0, one more step for each 9 in between), and once known it equals the true
carry. The done flag is $\mathbb 1[c_i\ne ?]$ and the committed output is
$(a_i+b_i+c_i)\bmod 10$, which is correct. The input halts at $\max_i\rho_i = d^*(x) = \mathrm{LB}(x)$
for $r=1$. \emph{Parity}: the same construction with the segment prefix parity in place of the carry,
initialised to $0$ at delimiters.
\end{proof}
